\documentclass[11pt]{amsart}

% Encoding and fonts
\usepackage[T1]{fontenc}
\usepackage[utf8]{inputenc}
\usepackage{lmodern}

% Math packages
\usepackage{amsmath}
\usepackage{amssymb}
\usepackage{amsfonts}
\usepackage{amsthm}
\usepackage{mathtools}

% Layout and utilities
\usepackage{geometry}
\usepackage{enumitem}

\geometry{margin=1in}

\numberwithin{equation}{section}
\setcounter{tocdepth}{3}
\setcounter{secnumdepth}{3}

% ---------------------------------------------------------
% Theorem-like environments
% ---------------------------------------------------------

\theoremstyle{plain}
\newtheorem{theorem}{Theorem}[section]
\newtheorem{proposition}[theorem]{Proposition}
\newtheorem{lemma}[theorem]{Lemma}
\newtheorem{corollary}[theorem]{Corollary}
\newtheorem{conjecture}[theorem]{Conjecture}

\theoremstyle{definition}
\newtheorem{definition}[theorem]{Definition}
\newtheorem{experiment}[theorem]{Experiment}
\newtheorem{heuristic}[theorem]{Heuristic}
\newtheorem{question}[theorem]{Question}

\theoremstyle{remark}
\newtheorem{remark}[theorem]{Remark}
\newtheorem{warning}[theorem]{Warning}

\newcommand{\R}{\mathbb{R}}
\newcommand{\C}{\mathbb{C}}
\newcommand{\N}{\mathbb{N}}
\newcommand{\Qp}{\mathbb{Q}_{>0}}
\newcommand{\Hil}{\mathcal{H}}
\newcommand{\Sch}{\mathcal{S}}
\newcommand{\Four}{\mathcal{F}}
\newcommand{\Dom}{\mathcal{D}}
\newcommand{\Bop}{\mathcal{B}}
\newcommand{\id}{\mathrm{Id}}
\newcommand{\ad}{\operatorname{ad}}
\newcommand{\Reop}{\operatorname{Re}}
\newcommand{\Om}{\Omega}
\newcommand{\twconv}{\star_{\vartheta}}
\newcommand{\deltaone}{\delta_{1}}
\newcommand{\ellft}{\lambda}

% hyperref must be loaded last
\usepackage[hidelinks]{hyperref}

\title[Heisenberg Uncertainty and M\"obius Dynamics IV]{
Heisenberg Uncertainty and Twisted M\"obius Dynamics IV:\
Cohomological Curvature, Prime--Factor Variance,\\
and a Zeta--Zero Singularity Criterion
}

\author{}
\date{}

\begin{document}

\maketitle

\tableofcontents

\bigskip

\begin{remark}[Version 4]
\label{rem:status}
The canonical pair, Weyl calculus, Robertson inequality, and general twisted
algebra are standard. The arithmetic specialization
\(n\mapsto(\ell\log n,\kappa\Omega(n))\) produces the specific cocycle studied
here. Version 4 adds three exact results: a structural obstruction showing
that every expectation value of a linearly M\"obius-weighted state is
necessarily quadratic in the M\"obius coefficients; a closed formula for the
second \(\vartheta\)-curvature of the twisted inverse \(z_\vartheta(n)\); and
an exact Dirichlet-series summation of that curvature which continues locally
to \(\Re s>1/2\) and develops a logarithmic branch singularity at every
Riemann zero in that half-plane.
\end{remark}

\section{Introduction: the Version 4 structure}
\label{sec:introduction}

Version 1 placed a standard Heisenberg pair next to a zeta multiplier but did
not couple the two structures.  Version 2 corrected that defect: the exact
canonical commutation relation deformed the arithmetic product itself through
a nontrivial cocycle.  The Version 2 audit nevertheless identified two
remaining weaknesses.  First, the actual uncertainty principle was not stated
prominently.  Second, several domain, inverse, and convergence claims required
repair.

Version 4 starts from the self-adjoint position and momentum operators
\begin{equation}
(Q\psi)(q)=q\psi(q),
\qquad
(P\psi)(q)=-i\hbar\psi'(q),
\label{eq:intro_QP}
\end{equation}
with the canonical commutation relation
\begin{equation}
[Q,P]=i\hbar I
\label{eq:intro_CCR}
\end{equation}
on the Schwartz core.  For every normalized Schwartz state, the Heisenberg uncertainty principle is
\begin{equation}
\boxed{\Delta_\psi Q\,\Delta_\psi P\ge\frac{\hbar}{2}.}
\label{eq:intro_heisenberg_uncertainty}
\end{equation}
Here the uncertainty is a mathematical spread intrinsic to the state, not an
experimental error.

The arithmetic specialization assigns to an integer \(n\) the phase-space
increments
\begin{equation}
x_n=\ell\log n,
\qquad
y_n=\kappa\Omega(n),
\label{eq:intro_phase_space_coordinates}
\end{equation}
and the self-adjoint generator
\begin{equation}
G_n=y_nQ-x_nP.
\label{eq:intro_Gn}
\end{equation}
The first central statement of Version 4 is an uncertainty law for two
integers:
\begin{equation}
\boxed{
\Delta_\psi G_m\,\Delta_\psi G_n
\ge
\frac{\hbar\ell\kappa}{2}
\left|\Omega(m)\log n-\log m\,\Omega(n)\right|.}
\label{eq:intro_arithmetic_uncertainty}
\end{equation}
For distinct primes \(p,q\), this becomes
\begin{equation}
\boxed{
\Delta_\psi G_p\,\Delta_\psi G_q
\ge
\frac{\hbar\ell\kappa}{2}
\left|\log\frac{q}{p}\right|.}
\label{eq:intro_prime_uncertainty}
\end{equation}
Thus the same arithmetic symplectic area that later appears in the twisted
M\"obius product already appears as a lower bound forced directly by the
Heisenberg uncertainty principle.

The associated Weyl operators are
\begin{equation}
W_n=\exp\left(\frac{i}{\hbar}G_n\right),
\label{eq:intro_Wn}
\end{equation}
and satisfy
\begin{equation}
W_mW_n
=
\chi_\vartheta(m,n)W_{mn},
\label{eq:intro_twisted_product}
\end{equation}
where
\begin{equation}
\chi_\vartheta(m,n)
=
\exp\left[
\frac{i\vartheta}{2}
\bigl(\Omega(m)\log n-\log m\,\Omega(n)\bigr)
\right],
\qquad
\vartheta=\frac{\ell\kappa}{\hbar}.
\label{eq:intro_cocycle}
\end{equation}
The cocycle is proved below to be cohomologically nontrivial for every
\(\vartheta\ne0\).  It induces an associative noncommutative Dirichlet
product.  M\"obius inversion in that product defines coefficients
\(z_\vartheta(n)\).  For prime powers and distinct prime pairs one obtains
\begin{align}
z_\vartheta(p^r)&=1,
\label{eq:intro_prime_power_formula}\\
z_\vartheta(pq)
&=2\cos\left(\frac{\vartheta}{2}\log\frac{q}{p}\right)-1.
\label{eq:intro_prime_pair_formula}
\end{align}
The same determinant-like arithmetic quantity therefore controls both the
uncertainty lower bound \eqref{eq:intro_arithmetic_uncertainty} and the phase
responsible for the deformed inverse.


Version 4 also differentiates the twisted inverse itself at the commuting
point. For
\begin{equation}
n=\prod_{j=1}^{r}p_j^{a_j},
\qquad
T(a)=\frac{a(a+1)}2,
\label{eq:intro_factorization_T}
\end{equation}
one obtains
\begin{equation}
\boxed{
-z_0''(n)=\frac12\sum_{1\le i<j\le r}
T(a_i)T(a_j)\log^2\!\frac{p_j}{p_i}.}
\label{eq:intro_curvature_formula}
\end{equation}
For square-free \(n=p_1\cdots p_r\),
\begin{equation}
\boxed{
-z_0''(n)=\frac12\left[r\sum_{p\mid n}(\log p)^2-(\log n)^2\right]
=\frac{r^2}{2}\operatorname{Var}_{p\mid n}(\log p).}
\label{eq:intro_variance_formula}
\end{equation}
Summing this curvature does not return a two-variable M\"obius correlation.
For \(\Re s>1\),
\begin{equation}
\mathcal K(s)=-\sum_{n\ge1}\frac{z_0''(n)}{n^s}
=\frac{\zeta(s)}2\bigl(S_0(s)S_2(s)-S_1(s)^2\bigr),
\label{eq:intro_curvature_determinant}
\end{equation}
where
\begin{equation}
S_j(s)=\sum_p\frac{p^{-s}}{(1-p^{-s})^2}(\log p)^j.
\label{eq:intro_Sj}
\end{equation}
Section~\ref{sec:curvature_dirichlet_series} rewrites these functions through
\(\log\zeta(rs)\) and \(\zeta'(rs)/\zeta(rs)\), producing a local
continuation to \(\Re s>1/2\) and an exact singularity criterion for
nontrivial zeta zeros to the right of the critical line.

Finally, a bounded self-adjoint M\"obius Hamiltonian is defined for
\(\sigma>1\).  Its Schr\"odinger evolution and the exact Heisenberg equations
for \(Q(t)\) and \(P(t)\) are derived with domains stated explicitly.  The
paper does not claim a connection to RH; analytic continuation, functional
equations, and spectral information remain open.

\begin{remark}[Scope of Version 4]
\label{rem:v3_scope}
The operator deformation is developed in the half-planes stated for each
operator. Version 4 additionally derives a direct local relation between the
second deformation curvature and the classical zero set in
Section~\ref{sec:curvature_zero_detection}.
\end{remark}

\section{Exact canonical kinematics}
\label{sec:kinematics}

\subsection{Hilbert space, position, and momentum}
\label{subsec:QP_definitions}

\begin{definition}[Hilbert space]
\label{def:hilbert_space}
Fix $\hbar>0$ and set
\begin{equation}
\Hil=L^2(\R,dq)
\label{eq:hilbert_space}
\end{equation}
with its standard complex Hilbert-space inner product.
\end{definition}

\begin{definition}[Position]
\label{def:position}
The position operator $Q$ is
\begin{equation}
(Q\psi)(q)=q\psi(q)
\label{eq:Q_action}
\end{equation}
with maximal domain
\begin{equation}
\Dom(Q)=\{\psi\in L^2(\R):q\psi(q)\in L^2(\R)\}.
\label{eq:Q_domain}
\end{equation}
\end{definition}

\begin{definition}[Momentum]
\label{def:momentum}
The momentum operator $P$ is
\begin{equation}
(P\psi)(q)=-i\hbar\,\psi'(q)
\label{eq:P_action}
\end{equation}
with domain
\begin{equation}
\Dom(P)=H^1(\R).
\label{eq:P_domain}
\end{equation}
\end{definition}

\begin{proposition}[Self-adjointness and common core]
\label{prop:self_adjoint_QP}
The operators $Q$ and $P$ are self-adjoint.  The Schwartz space
$\Sch(\R)$ is a common invariant core for both.
\end{proposition}

\begin{proof}
$Q$ is the maximal multiplication operator by the real function $q$.  Under
the unitary Fourier transform with convention
\begin{equation}
(\Four\psi)(p)
=
\frac{1}{\sqrt{2\pi\hbar}}
\int_{\R}e^{-ipq/\hbar}\psi(q)\,dq,
\label{eq:fourier_transform}
\end{equation}
$P$ is unitarily equivalent to multiplication by the real variable $p$.
The core statement follows by standard cutoff and mollification arguments.
See, for example, \cite{ReedSimon1980}.
\end{proof}

\begin{proposition}[Exact canonical commutation relation]
\label{prop:CCR}
For every $\psi\in\Sch(\R)$,
\begin{equation}
[Q,P]\psi=i\hbar\psi.
\label{eq:CCR}
\end{equation}
\end{proposition}

\begin{proof}
Direct differentiation gives
\begin{align}
(QP\psi)(q)
&=-i\hbar q\psi'(q),
\label{eq:QP}\\
(PQ\psi)(q)
&=-i\hbar\psi(q)-i\hbar q\psi'(q).
\label{eq:PQ}
\end{align}
Subtracting \eqref{eq:PQ} from \eqref{eq:QP} proves \eqref{eq:CCR}.
\end{proof}

\subsection{The Heisenberg uncertainty principle and exact incompatibility}
\label{subsec:heisenberg_uncertainty}

\begin{definition}[Expectation and uncertainty]
\label{def:uncertainty}
Let \(A\) be self-adjoint and let \(\psi\in\Dom(A)\) be normalized.  Define
\begin{align}
\langle A\rangle_\psi
&=\langle\psi,A\psi\rangle,
\label{eq:expectation_definition}\\
(\Delta_\psi A)^2
&=\|(A-\langle A\rangle_\psi)\psi\|^2.
\label{eq:variance_definition}
\end{align}
\end{definition}

\begin{theorem}[Heisenberg uncertainty principle]
\label{thm:heisenberg_uncertainty}
For every normalized \(\psi\in\Sch(\R)\),
\begin{equation}
\boxed{
\Delta_\psi Q\,\Delta_\psi P
\ge
\frac{\hbar}{2}.}
\label{eq:heisenberg_uncertainty}
\end{equation}
\end{theorem}

\begin{proof}
Set
\begin{equation}
\widetilde Q=Q-\langle Q\rangle_\psi I,
\qquad
\widetilde P=P-\langle P\rangle_\psi I.
\label{eq:centered_QP}
\end{equation}
By Cauchy--Schwarz,
\begin{equation}
\Delta_\psi Q\,\Delta_\psi P
\ge
\left|\langle\widetilde Q\psi,\widetilde P\psi\rangle\right|.
\label{eq:CS_uncertainty}
\end{equation}
The modulus dominates the absolute value of the imaginary part, and on
\(\Sch(\R)\)
\begin{align}
2i\,\operatorname{Im}
\langle\widetilde Q\psi,\widetilde P\psi\rangle
&=
\langle\psi,[Q,P]\psi\rangle
\label{eq:imaginary_commutator}\\
&=i\hbar.
\label{eq:imaginary_CCR}
\end{align}
Combining \eqref{eq:CS_uncertainty}--\eqref{eq:imaginary_CCR} proves
\eqref{eq:heisenberg_uncertainty}.
\end{proof}

\begin{remark}[Meaning of the inequality]
\label{rem:uncertainty_meaning}
Equation~\eqref{eq:heisenberg_uncertainty} is the Heisenberg uncertainty
principle used in this paper.  The quantities \(\Delta_\psi Q\) and
\(\Delta_\psi P\) are intrinsic standard deviations of a state.  They are
not measurement errors.  Exact simultaneous localization would require both
to vanish, which is prohibited by \eqref{eq:heisenberg_uncertainty}.
\end{remark}

\begin{theorem}[No simultaneous exact position and momentum]
\label{thm:no_joint_eigenstate}
There is no nonzero state with simultaneous exact position and exact momentum.
More precisely, both of the following are impossible for real \(q_0,p_0\):
\begin{enumerate}[leftmargin=*]
\item a nonzero vector \(\psi\) in a common invariant domain satisfying
\begin{equation}
Q\psi=q_0\psi,
\qquad
P\psi=p_0\psi;
\label{eq:joint_eigenstate_assumption}
\end{equation}
\item a nonzero tempered distribution \(T\in\Sch'(\R)\) satisfying
\begin{equation}
qT=q_0T,
\qquad
-i\hbar T'=p_0T.
\label{eq:joint_distribution_assumption}
\end{equation}
\end{enumerate}
\end{theorem}

\begin{proof}
For vectors, \eqref{eq:joint_eigenstate_assumption} gives
\([Q,P]\psi=0\), whereas the CCR gives \([Q,P]\psi=i\hbar\psi\); hence
\(\psi=0\).  For distributions define directly
\begin{equation}
QT=qT,
\qquad
PT=-i\hbar T'.
\label{eq:distribution_QP_definition}
\end{equation}
The distributional Leibniz rule gives
\begin{equation}
[Q,P]T=i\hbar T.
\label{eq:distribution_CCR}
\end{equation}
The simultaneous eigenvalue equations would make the left-hand side zero,
so \(T=0\).
\end{proof}

\section{Weyl displacements and the arithmetic phase space}
\label{sec:weyl_arithmetic}

\subsection{Weyl operators}
\label{subsec:weyl_operators}

\begin{definition}[Weyl displacement]
\label{def:weyl_operator}
For $(x,y)\in\R^2$, define
\begin{equation}
W(x,y)
=
\exp\!\left(\frac{i}{\hbar}(yQ-xP)\right).
\label{eq:weyl_definition}
\end{equation}
This is a unitary operator on all of $\Hil$.
\end{definition}

\begin{proposition}[Coordinate action]
\label{prop:weyl_coordinate_action}
For $\psi\in L^2(\R)$,
\begin{equation}
(W(x,y)\psi)(q)
=
\exp\!\left(\frac{i}{\hbar}y\left(q-\frac{x}{2}\right)\right)
\psi(q-x)
\label{eq:weyl_coordinate_action}
\end{equation}
almost everywhere.
\end{proposition}

\begin{proof}
Since $[iyQ/\hbar,-ixP/\hbar]$ is scalar, the Baker--Campbell--Hausdorff
formula gives
\begin{equation}
W(x,y)
=
e^{-ixy/(2\hbar)}e^{iyQ/\hbar}e^{-ixP/\hbar}.
\label{eq:weyl_BCH_factorization}
\end{equation}
The last factor translates by $x$ and the middle factor multiplies by
$e^{iyq/\hbar}$, yielding \eqref{eq:weyl_coordinate_action}.
\end{proof}

\begin{proposition}[Weyl form of the CCR]
\label{prop:weyl_relation}
Let $z=(x,y)$ and $z'=(x',y')$.  Define
\begin{equation}
\omega(z,z')=yx'-xy'.
\label{eq:symplectic_form}
\end{equation}
Then
\begin{equation}
W(z)W(z')
=
\exp\!\left(\frac{i}{2\hbar}\omega(z,z')\right)
W(z+z').
\label{eq:weyl_product}
\end{equation}
Consequently,
\begin{equation}
[W(z),W(z')]
=
2i\sin\!\left(\frac{\omega(z,z')}{2\hbar}\right)W(z+z').
\label{eq:weyl_commutator}
\end{equation}
\end{proposition}

\begin{proof}
For
\begin{equation}
A_z=\frac{i}{\hbar}(yQ-xP),
\label{eq:Az_definition}
\end{equation}
the CCR gives
\begin{equation}
[A_z,A_{z'}]
=
\frac{i}{\hbar}\omega(z,z')I.
\label{eq:Az_commutator}
\end{equation}
This commutator is central, so the Baker--Campbell--Hausdorff series stops
after the commutator term and gives \eqref{eq:weyl_product}.  Reversing $z$
and $z'$ changes the sign of $\omega$, giving \eqref{eq:weyl_commutator}.
The Weyl relations are standard; see \cite{Folland1989}.
\end{proof}

\begin{proposition}[Position and momentum displacement equations]
\label{prop:weyl_QP_commutators}
On $\Sch(\R)$,
\begin{align}
[Q,W(x,y)]
&=xW(x,y),
\label{eq:Q_W_commutator}\\
[P,W(x,y)]
&=yW(x,y).
\label{eq:P_W_commutator}
\end{align}
\end{proposition}

\begin{proof}
Both identities follow by differentiating \eqref{eq:weyl_coordinate_action},
or directly from \eqref{eq:CCR} because the relevant commutators with the
exponent are scalar.
\end{proof}

\subsection{Arithmetic embedding}
\label{subsec:arithmetic_embedding}

\begin{definition}[The additive prime-factor counter]
\label{def:Omega}
For
\begin{equation}
n=\prod_p p^{v_p(n)},
\label{eq:prime_factorization}
\end{equation}
define
\begin{equation}
\Om(n)=\sum_p v_p(n).
\label{eq:Omega_definition}
\end{equation}
Then
\begin{equation}
\Om(mn)=\Om(m)+\Om(n).
\label{eq:Omega_additivity}
\end{equation}
\end{definition}

\begin{definition}[Arithmetic phase-space map]
\label{def:arithmetic_phase_space_map}
Fix two positive scales $\ell>0$ and $\kappa>0$.  For $n\in\N$, put
\begin{align}
x_n
&=\ell\log n,
\label{eq:xn_definition}\\
y_n
&=\kappa\Om(n),
\label{eq:yn_definition}\\
z_n
&=(x_n,y_n).
\label{eq:zn_definition}
\end{align}
Define the dimensionless Heisenberg area parameter
\begin{equation}
\vartheta=\frac{\ell\kappa}{\hbar}.
\label{eq:theta_definition}
\end{equation}
\end{definition}

\begin{remark}[Why these two coordinates]
\label{rem:why_coordinates}
The first coordinate is forced by multiplicative size because
$\log(mn)=\log m+\log n$.  The second coordinate is chosen as the simplest
integer-valued completely additive function that directly controls the
M\"obius sign on square-free integers:
\begin{equation}
\mu(n)\neq0
\quad\Longrightarrow\quad
\mu(n)=(-1)^{\Om(n)}.
\label{eq:mobius_Omega_relation}
\end{equation}
The construction can be generalized by replacing $\Om$ with another
completely additive arithmetic function, but the present notebook fixes $\Om$ throughout.
\end{remark}

\begin{definition}[Arithmetic Weyl operators]
\label{def:arithmetic_weyl_operator}
For $n\in\N$, define
\begin{equation}
W_n=W(x_n,y_n).
\label{eq:Wn_definition}
\end{equation}
\end{definition}

\begin{theorem}[Exact Heisenberg cocycle on the integers]
\label{thm:integer_weyl_cocycle}
For $m,n\in\N$,
\begin{equation}
W_mW_n
=
\chi_{\vartheta}(m,n)W_{mn},
\label{eq:integer_weyl_product}
\end{equation}
where
\begin{equation}
\chi_{\vartheta}(m,n)
=
\exp\!\left[
\frac{i\vartheta}{2}
\bigl(\Om(m)\log n-\log m\,\Om(n)\bigr)
\right].
\label{eq:chi_definition}
\end{equation}
Moreover,
\begin{equation}
[W_m,W_n]
=
2i\sin\!\left[
\frac{\vartheta}{2}
\bigl(\Om(m)\log n-\log m\,\Om(n)\bigr)
\right]W_{mn}.
\label{eq:integer_weyl_commutator}
\end{equation}
\end{theorem}

\begin{proof}
Complete additivity gives
\begin{equation}
z_{mn}=z_m+z_n.
\label{eq:phase_space_additivity}
\end{equation}
Substitution of \eqref{eq:xn_definition}--\eqref{eq:theta_definition} into
\eqref{eq:weyl_product} gives \eqref{eq:integer_weyl_product};
\eqref{eq:integer_weyl_commutator} follows from
\eqref{eq:weyl_commutator}.
\end{proof}

\begin{corollary}[Prime commutator]
\label{cor:prime_commutator}
For distinct primes $p,q$,
\begin{equation}
[W_p,W_q]
=
2i\sin\!\left(\frac{\vartheta}{2}\log\frac{q}{p}\right)W_{pq}.
\label{eq:prime_commutator}
\end{equation}
Thus the prime operators commute only at the resonant values
\begin{equation}
\vartheta\log\frac{q}{p}\in2\pi\mathbb Z.
\label{eq:prime_commutation_resonance}
\end{equation}
\end{corollary}

\begin{remark}[Where Heisenberg enters nontrivially]
\label{rem:where_heisenberg_enters}
Equation~\eqref{eq:integer_weyl_product} is the first point where Version 2
ceases to be ordinary Dirichlet algebra in operator notation.  If the
canonical pair were commuting, the Weyl cocycle would be $1$ and the
arithmetic product below would reduce to ordinary Dirichlet convolution.
For nonzero $\vartheta$, the CCR creates order-dependent phases between
different arithmetic factors.
\end{remark}

\begin{theorem}[The arithmetic cocycle is not gauge-trivial]
\label{thm:cocycle_nontrivial}
Extend \(\log\) and \(\Omega\) additively to
\(\mathbb Q_{>0}^{\times}\).  For every real \(\vartheta\ne0\), the
cocycle \(\chi_\vartheta\) is not a scalar coboundary on
\(\mathbb N^{\times}\), and its extension is not a scalar coboundary on
\(\mathbb Q_{>0}^{\times}\).
\end{theorem}

\begin{proof}
On a commutative semigroup, every scalar coboundary
\begin{equation}
(\delta a)(m,n)=a(m)a(n)a(mn)^{-1}
\label{eq:scalar_coboundary}
\end{equation}
is symmetric in \(m,n\).  Hence a necessary condition for
\(\chi_\vartheta\) to be a coboundary is
\begin{equation}
\frac{\chi_\vartheta(m,n)}{\chi_\vartheta(n,m)}=1
\label{eq:cocycle_commutator_trivial}
\end{equation}
for all \(m,n\).  For distinct primes,
\begin{equation}
\frac{\chi_\vartheta(p,q)}{\chi_\vartheta(q,p)}
=
\exp\left(i\vartheta\log\frac{q}{p}\right).
\label{eq:prime_cocycle_commutator}
\end{equation}
If this were one simultaneously for \((p,q)=(2,3)\) and \((2,5)\), then
there would be nonzero integers \(k,l\) such that
\begin{equation}
\vartheta\log(3/2)=2\pi k,
\qquad
\vartheta\log(5/2)=2\pi l.
\label{eq:two_prime_phase_integrality}
\end{equation}
Thus \(\log(3/2)/\log(5/2)=k/l\in\mathbb Q\).  Exponentiating a
reduced rational relation would produce an equality between incompatible
prime factorizations, contradicting unique factorization.  Therefore
\eqref{eq:cocycle_commutator_trivial} cannot hold for all primes.
\end{proof}


\section{Structural filter: linear M\"obius states cannot escape second order}
\label{sec:linear_state_filter}

Let \(v_n\in\Hil\) be arbitrary vectors and define
\begin{equation}
\Psi_x=\sum_{n\le x}\mu(n)v_n.
\label{eq:linear_mobius_state}
\end{equation}

\begin{theorem}[Quadratic obstruction for all expectation-value routes]
\label{thm:linear_state_quadratic_obstruction}
For every linear operator \(A\) for which the expression is defined,
\begin{equation}
\langle\Psi_x,A\Psi_x\rangle
=\sum_{m,n\le x}\mu(m)\mu(n)\langle v_m,Av_n\rangle.
\label{eq:expectation_bilinear_mobius}
\end{equation}
Hence every norm, expectation value, covariance, variance, or Robertson bound
obtained from a state constructed linearly with coefficients \(\mu(n)\) is a
quadratic M\"obius form.
\end{theorem}

\begin{proof}
Expand the inner product sesquilinearly and use that \(\mu(n)\) is real.
\end{proof}

\begin{corollary}[State-superposition route closed]
\label{cor:linear_state_route_closed}
Changing the reference vector, the vectors \(v_n\), or the observable \(A\)
does not alter the quadratic dependence in
\eqref{eq:expectation_bilinear_mobius}. A non-quadratic consequence must
therefore be extracted from the algebraic deformation itself.
\end{corollary}

\section{Arithmetic uncertainty generated by the CCR}
\label{sec:arithmetic_uncertainty}

\begin{definition}[Integer generator]
\label{def:integer_generator}
For \(n\ge1\), define on \(\Sch(\R)\) the symmetric operator
\begin{equation}
G_n^{(0)}=y_nQ-x_nP
=\kappa\Omega(n)Q-\ell\log n\,P.
\label{eq:Gn_definition}
\end{equation}
It is essentially self-adjoint on \(\Sch(\R)\); denote its self-adjoint
closure by \(G_n\).  The corresponding unitary displacement is
\begin{equation}
W_n=\exp(iG_n/\hbar).
\label{eq:Wn_from_Gn}
\end{equation}
\end{definition}

\begin{proposition}[Arithmetic commutator]
\label{prop:arithmetic_generator_commutator}
For all \(m,n\ge1\), the restrictions of \(G_m\) and \(G_n\) to \(\Sch(\R)\) satisfy
\begin{align}
[G_m,G_n]
&=i\hbar(x_my_n-y_mx_n)I
\label{eq:GmGn_commutator_xy}\\
&=i\hbar\ell\kappa
\bigl(\log m\,\Omega(n)-\Omega(m)\log n\bigr)I.
\label{eq:GmGn_commutator_arithmetic}
\end{align}
\end{proposition}

\begin{proof}
Expand the commutator and use \([Q,P]=i\hbar I\).  Terms involving
\([Q,Q]\) and \([P,P]\) vanish.
\end{proof}

\begin{theorem}[Arithmetic Heisenberg uncertainty law]
\label{thm:arithmetic_uncertainty}
For every normalized \(\psi\in\Sch(\R)\) and all \(m,n\ge1\),
\begin{equation}
\boxed{
\Delta_\psi G_m\,\Delta_\psi G_n
\ge
\frac{\hbar\ell\kappa}{2}
\left|\Omega(m)\log n-\log m\,\Omega(n)\right|.}
\label{eq:arithmetic_uncertainty}
\end{equation}
For distinct primes \(p,q\),
\begin{equation}
\boxed{
\Delta_\psi G_p\,\Delta_\psi G_q
\ge
\frac{\hbar\ell\kappa}{2}
\left|\log\frac{q}{p}\right|.}
\label{eq:prime_pair_uncertainty}
\end{equation}
For two powers of the same prime the lower bound is zero.
\end{theorem}

\begin{proof}
Apply the Robertson inequality to the self-adjoint pair \(G_m,G_n\):
\begin{equation}
\Delta_\psi G_m\,\Delta_\psi G_n
\ge
\frac12
\left|\langle\psi,[G_m,G_n]\psi\rangle\right|.
\label{eq:Robertson_GmGn}
\end{equation}
Substitution of \eqref{eq:GmGn_commutator_arithmetic} gives
\eqref{eq:arithmetic_uncertainty}.  If \(p,q\) are distinct primes then
\(\Omega(p)=\Omega(q)=1\), which gives \eqref{eq:prime_pair_uncertainty}.
If \(m=p^a,n=p^b\), the determinant
\(\Omega(m)\log n-\log m\,\Omega(n)\) vanishes.
\end{proof}

\begin{remark}[The common invariant]
\label{rem:common_symplectic_invariant}
Define the arithmetic symplectic area
\begin{equation}
\mathcal A(m,n)
=
\Omega(m)\log n-\log m\,\Omega(n).
\label{eq:arithmetic_symplectic_area}
\end{equation}
The uncertainty lower bound is proportional to \( |\mathcal A(m,n)| \),
while the twisted product below carries the phase
\(\exp(i\vartheta\mathcal A(m,n)/2)\).  Thus uncertainty and twisted
M\"obius inversion are controlled by the same arithmetic invariant, although
they are distinct consequences of the CCR.
\end{remark}

\subsection{The semiclassical and arithmetic-commutative limits are distinct}
\label{subsec:two_limits}

The parameter controlling the Weyl cocycle is
\begin{equation}
\vartheta(\hbar,\ell,\kappa)
=
\frac{\ell\kappa}{\hbar}.
\label{eq:vartheta_scaling}
\end{equation}
The uncertainty scale and the deformation scale therefore depend on the same
three parameters in different ways.

\begin{proposition}[Opposite fixed-scale limits]
\label{prop:opposite_limits}
Fix \(\ell>0\) and \(\kappa>0\).  For every fixed pair \(m,n\),
\begin{equation}
\frac{\hbar\ell\kappa}{2}|\mathcal A(m,n)|
\longrightarrow0
\qquad(\hbar\to0^+),
\label{eq:uncertainty_classical_limit}
\end{equation}
whereas
\begin{equation}
|\vartheta|
=
\frac{\ell\kappa}{\hbar}
\longrightarrow\infty.
\label{eq:twist_fixed_scale_limit}
\end{equation}
Consequently the ordinary semiclassical limit \(\hbar\to0^+\) with
\(\ell,\kappa\) fixed does not approach the untwisted Dirichlet product.
\end{proposition}

\begin{proof}
Equation~\eqref{eq:uncertainty_classical_limit} follows directly from
Theorem~\ref{thm:arithmetic_uncertainty}, while
\eqref{eq:twist_fixed_scale_limit} follows from
\eqref{eq:vartheta_scaling}.
\end{proof}

\begin{proposition}[Joint semiclassical--commutative scaling]
\label{prop:joint_scaling}
Let \(\ell=\ell(\hbar)>0\) and
\(\kappa=\kappa(\hbar)>0\).  Then
\begin{equation}
\vartheta(\hbar)\longrightarrow0
\qquad(\hbar\to0^+)
\label{eq:theta_to_zero_condition}
\end{equation}
if and only if
\begin{equation}
\ell(\hbar)\kappa(\hbar)=o(\hbar).
\label{eq:joint_scaling_condition}
\end{equation}
Under this scaling the twisted product converges coefficientwise to the
ordinary Dirichlet product, and the arithmetic uncertainty lower bound obeys
\begin{equation}
\frac{\hbar\ell(\hbar)\kappa(\hbar)}{2}
|\mathcal A(m,n)|
=o(\hbar^2)
\label{eq:joint_uncertainty_decay}
\end{equation}
for every fixed \(m,n\).
\end{proposition}

\begin{proof}
The equivalence of \eqref{eq:theta_to_zero_condition} and
\eqref{eq:joint_scaling_condition} is immediate from
\eqref{eq:vartheta_scaling}.  For fixed \(m,n\),
\(\chi_{\vartheta}(m,n)\to1\) as \(\vartheta\to0\), so the
finite divisor sum defining \(f\star_{\vartheta}g\) converges to ordinary
Dirichlet convolution coefficientwise.  Multiplying
\(\ell(\hbar)\kappa(\hbar)=o(\hbar)\) by \(\hbar/2\) gives
\eqref{eq:joint_uncertainty_decay}.
\end{proof}

\begin{remark}[Two independent notions of classicality]
\label{rem:two_classicalities}
The limit \(\hbar\to0^+\) suppresses the Robertson--Heisenberg lower
bound, whereas \(\vartheta\to0\) removes the arithmetic twist.  They are
not the same operation.  With fixed \(\ell,\kappa\) they move in opposite
directions; simultaneous suppression requires the joint scaling
\eqref{eq:joint_scaling_condition}.
\end{remark}

\section{The Heisenberg-twisted Dirichlet algebra}
\label{sec:twisted_dirichlet_algebra}

\subsection{Definition and associativity}
\label{subsec:twisted_convolution_definition}

\begin{definition}[Heisenberg-twisted Dirichlet convolution]
\label{def:twisted_convolution}
For arithmetic functions $f,g:\N\to\C$, define
\begin{equation}
(f\twconv g)(k)
=
\sum_{d\mid k}
f(d)g(k/d)\chi_{\vartheta}(d,k/d).
\label{eq:twisted_convolution}
\end{equation}
The sum is finite for every $k$.
\end{definition}

\begin{lemma}[Cocycle identity]
\label{lem:cocycle_identity}
For all $a,b,c\in\N$,
\begin{equation}
\chi_{\vartheta}(a,b)\chi_{\vartheta}(ab,c)
=
\chi_{\vartheta}(b,c)\chi_{\vartheta}(a,bc).
\label{eq:cocycle_identity}
\end{equation}
\end{lemma}

\begin{proof}
The exponent in \eqref{eq:chi_definition} is the alternating bilinear form
obtained from the additive map $n\mapsto z_n$.  Expanding both sides using
$\log(ab)=\log a+\log b$ and
$\Om(ab)=\Om(a)+\Om(b)$ gives the same exponent.
\end{proof}

\begin{proposition}[Associative noncommutative arithmetic algebra]
\label{prop:twisted_associative}
The product $\twconv$ is associative and has identity $\deltaone$, where
\begin{equation}
\deltaone(1)=1,
\qquad
\deltaone(n)=0\quad(n>1).
\label{eq:delta_one_definition}
\end{equation}
For $\vartheta\neq0$ it is generally noncommutative.
\end{proposition}

\begin{proof}
Associativity follows from Lemma~\ref{lem:cocycle_identity} after grouping
the finite divisor sums by triples $abc=k$.  The identity property follows
from $\chi_{\vartheta}(1,n)=\chi_{\vartheta}(n,1)=1$.  Noncommutativity
follows from Corollary~\ref{cor:prime_commutator}; equivalently,
\begin{equation}
(\delta_p\twconv\delta_q)(pq)
=
\chi_{\vartheta}(p,q)
\label{eq:delta_pq_forward}
\end{equation}
whereas
\begin{equation}
(\delta_q\twconv\delta_p)(pq)
=
\chi_{\vartheta}(q,p)
=
\overline{\chi_{\vartheta}(p,q)}.
\label{eq:delta_pq_reverse}
\end{equation}
\end{proof}

\subsection{Weighted Banach algebra and operator representation}
\label{subsec:weighted_banach_algebra}

\begin{definition}[Weighted arithmetic Banach space]
\label{def:weighted_banach_algebra}
For $\sigma>0$, set
\begin{equation}
\mathcal A_{\sigma}
=
\left\{f:\N\to\C:
\|f\|_{\sigma}:=
\sum_{n\ge1}\frac{|f(n)|}{n^{\sigma}}<\infty
\right\}.
\label{eq:A_sigma_definition}
\end{equation}
\end{definition}

\begin{proposition}[Banach-algebra estimate]
\label{prop:banach_algebra_estimate}
For $f,g\in\mathcal A_{\sigma}$,
\begin{equation}
\|f\twconv g\|_{\sigma}
\le
\|f\|_{\sigma}\|g\|_{\sigma}.
\label{eq:banach_algebra_estimate}
\end{equation}
Hence $(\mathcal A_{\sigma},\twconv)$ is a unital Banach algebra.
\end{proposition}

\begin{proof}
Since $|\chi_{\vartheta}|=1$,
\begin{align}
\|f\twconv g\|_{\sigma}
&\le
\sum_{k\ge1}\frac{1}{k^{\sigma}}
\sum_{d\mid k}|f(d)||g(k/d)|
\label{eq:banach_estimate_step1}\\
&=
\sum_{d,e\ge1}
\frac{|f(d)|}{d^{\sigma}}
\frac{|g(e)|}{e^{\sigma}}
\label{eq:banach_estimate_step2}\\
&=
\|f\|_{\sigma}\|g\|_{\sigma}.
\label{eq:banach_estimate_step3}
\end{align}
Completeness is the completeness of weighted $\ell^1$.
\end{proof}

\begin{definition}[Weyl--Dirichlet representation]
\label{def:Weyl_Dirichlet_representation}
For $f\in\mathcal A_{\sigma}$, define
\begin{equation}
\Pi_{\sigma,\vartheta}(f)
=
\sum_{n\ge1}\frac{f(n)}{n^{\sigma}}W_n,
\label{eq:Pi_definition}
\end{equation}
where the series converges absolutely in operator norm.
\end{definition}

\begin{proposition}[Representation theorem]
\label{prop:representation_theorem}
For $f,g\in\mathcal A_{\sigma}$,
\begin{align}
\|\Pi_{\sigma,\vartheta}(f)\|
&\le\|f\|_{\sigma},
\label{eq:Pi_norm_bound}\\
\Pi_{\sigma,\vartheta}(f)
\Pi_{\sigma,\vartheta}(g)
&=
\Pi_{\sigma,\vartheta}(f\twconv g).
\label{eq:Pi_homomorphism}
\end{align}
\end{proposition}

\begin{proof}
The norm bound follows from $\|W_n\|=1$.  Absolute norm convergence permits
multiplication and regrouping of the double series.  The coefficient of
$W_k$ is exactly \eqref{eq:twisted_convolution} by
Theorem~\ref{thm:integer_weyl_cocycle}.
\end{proof}

\section{M\"obius inversion deformed by the CCR}
\label{sec:mobius_inversion}

\subsection{The M\"obius operator}
\label{subsec:mobius_operator}

\begin{definition}[M\"obius operator]
\label{def:mobius_operator}
For $\sigma>1$, define
\begin{equation}
\mathcal M_{\sigma,\vartheta}
=
\Pi_{\sigma,\vartheta}(\mu)
=
\sum_{n\ge1}
\frac{\mu(n)}{n^{\sigma}}W_n.
\label{eq:mobius_operator}
\end{equation}
The series converges absolutely in operator norm, so
$\mathcal M_{\sigma,\vartheta}\in\Bop(\Hil)$ and
\begin{equation}
\Dom(\mathcal M_{\sigma,\vartheta})=\Hil.
\label{eq:mobius_operator_domain}
\end{equation}
\end{definition}

\begin{remark}[Concrete coordinate action]
\label{rem:mobius_coordinate_action}
For $\psi\in L^2(\R)$,
\begin{equation}
(\mathcal M_{\sigma,\vartheta}\psi)(q)
=
\sum_{n\ge1}
\frac{\mu(n)}{n^{\sigma}}
\exp\!\left[
\frac{i}{\hbar}y_n\left(q-\frac{x_n}{2}\right)
\right]
\psi(q-x_n),
\label{eq:mobius_coordinate_action}
\end{equation}
with convergence in $L^2$.  Thus the operator is a M\"obius-weighted
superposition of simultaneous translations in position and momentum; it is
not a Fourier multiplier when $\kappa\neq0$.
\end{remark}

\subsection{Formal twisted inverse}
\label{subsec:formal_twisted_inverse}

\begin{proposition}[Existence and uniqueness of arithmetic inverses]
\label{prop:formal_inverse}
Let $f:\N\to\C$ satisfy $f(1)\neq0$.  Then there exists a unique arithmetic
function $g$ such that
\begin{equation}
f\twconv g=g\twconv f=\deltaone.
\label{eq:two_sided_inverse}
\end{equation}
Its coefficients are determined recursively from $g(1)=f(1)^{-1}$.
\end{proposition}

\begin{proof}
For $n>1$, the equation $(f\twconv r)(n)=0$ contains the term
$f(1)r(n)$ and otherwise only values $r(d)$ with $d<n$.  It therefore
constructs a unique right inverse $r$ satisfying $f\twconv r=\deltaone$.
Separately, the equations $(\ellft\twconv f)(n)=0$ construct a unique left
inverse $\ellft$ satisfying $\ellft\twconv f=\deltaone$.  Associativity
then gives
\begin{equation}
\ellft
=\ellft\twconv(f\twconv r)
=(\ellft\twconv f)\twconv r
=r,
\label{eq:left_right_inverse_coincide}
\end{equation}
so the inverse is bilateral and unique.
\end{proof}

\begin{definition}[Heisenberg-zeta coefficients]
\label{def:heisenberg_zeta_coefficients}
Let $z_{\vartheta}:\N\to\C$ be the unique two-sided inverse of $\mu$ for
$\twconv$:
\begin{equation}
\mu\twconv z_{\vartheta}
=
z_{\vartheta}\twconv\mu
=
\deltaone.
\label{eq:ztheta_inverse_definition}
\end{equation}
Thus
\begin{equation}
z_{\vartheta}(1)=1,
\label{eq:ztheta_one}
\end{equation}
and for $n>1$,
\begin{equation}
z_{\vartheta}(n)
=
-\sum_{\substack{d\mid n\\ d>1}}
\mu(d)z_{\vartheta}(n/d)
\chi_{\vartheta}(d,n/d).
\label{eq:ztheta_recursion}
\end{equation}
\end{definition}

\begin{proposition}[Commuting boundary]
\label{prop:commuting_boundary}
At $\vartheta=0$,
\begin{equation}
z_0(n)=1
\label{eq:z0_constant_one}
\end{equation}
for every $n\ge1$.
\end{proposition}

\begin{proof}
When $\vartheta=0$, $\chi_0\equiv1$ and $\twconv$ is ordinary Dirichlet
convolution.  Equation~\eqref{eq:ztheta_inverse_definition} becomes
$\mu*z_0=\deltaone$.  Classical M\"obius inversion gives $z_0=\mathbf1$.
See \cite{Apostol1976}.
\end{proof}

\begin{proposition}[Ordered-factorization formula]
\label{prop:ordered_factorization_formula}
For \(n>1\), the twisted inverse has the finite coefficient expansion
\begin{align}
z_\vartheta(n)
=
\sum_{r\ge1}(-1)^r
\sum_{\substack{n_1\cdots n_r=n\\ n_j>1}}
&\left(\prod_{j=1}^r\mu(n_j)\right)
\notag\\
&\times
\exp\left[
\frac{i\vartheta}{2}
\sum_{1\le i<j\le r}
\bigl(
\Omega(n_i)\log n_j-
\log n_i\,\Omega(n_j)
\bigr)
\right].
\label{eq:ordered_factorization_formula}
\end{align}
Only finitely many \(r\) and ordered factorizations contribute for fixed
\(n\).
\end{proposition}

\begin{proof}
Write \(\mu=\deltaone+h\) with \(h(1)=0\).  Coefficientwise the formal
Neumann expansion is
\begin{equation}
z_\vartheta
=
\sum_{r\ge0}(-h)^{\star_\vartheta r}.
\label{eq:formal_neumann_coefficientwise}
\end{equation}
For fixed \(n\), a product of more than \(\log_2 n\) factors exceeding
one cannot equal \(n\), so the coefficient sum is finite.  Iterating the
bicharacter cocycle yields the pairwise phase in
\eqref{eq:ordered_factorization_formula}.
\end{proof}

\subsection{A bounded inverse in a rigorous half-plane}
\label{subsec:bounded_inverse}

\begin{definition}[Uniform Neumann threshold]
\label{def:neumann_threshold}
Let \(\sigma_*\) be the unique real number \(>1\) satisfying
\begin{equation}
\frac{\zeta(\sigma_*)}{\zeta(2\sigma_*)}=2.
\label{eq:neumann_threshold_equation}
\end{equation}
Numerically,
\begin{equation}
\sigma_*=1.578017063034\ldots.
\label{eq:neumann_threshold_numeric}
\end{equation}
\end{definition}

\begin{theorem}[Uniform Neumann region]
\label{thm:neumann_region}
For every real \(\vartheta\) and every \(\sigma>\sigma_*\),
\begin{equation}
z_\vartheta\in\mathcal A_\sigma.
\label{eq:ztheta_in_Asigma}
\end{equation}
Consequently
\begin{equation}
\mathcal Z_{\sigma,\vartheta}
=
\Pi_{\sigma,\vartheta}(z_\vartheta)
=
\sum_{n\ge1}
\frac{z_\vartheta(n)}{n^\sigma}W_n
\label{eq:zeta_operator_deformed}
\end{equation}
converges absolutely in operator norm and satisfies
\begin{equation}
\mathcal M_{\sigma,\vartheta}
\mathcal Z_{\sigma,\vartheta}
=
\mathcal Z_{\sigma,\vartheta}
\mathcal M_{\sigma,\vartheta}
=I.
\label{eq:operator_inverse_identity}
\end{equation}
\end{theorem}

\begin{proof}
Write \(\mu=\deltaone+h\), where \(h(1)=0\).  For \(\sigma>1\),
\begin{align}
\|h\|_\sigma
&=
\sum_{n\ge2}\frac{\mu(n)^2}{n^\sigma}
\label{eq:h_norm_sigma_step1}\\
&=
\frac{\zeta(\sigma)}{\zeta(2\sigma)}-1.
\label{eq:h_norm_sigma_step2}
\end{align}
The right-hand side is continuous and strictly decreasing from \(+\infty\)
to \(0\) as \(\sigma\) runs from \(1^+\) to \(+\infty\), so
\(\sigma_*\) is unique and \(\|h\|_\sigma<1\) exactly when
\(\sigma>\sigma_*\).  The Banach-algebra Neumann series
\begin{equation}
(\deltaone+h)^{-1}
=
\sum_{r=0}^{\infty}(-h)^{\star_\vartheta r}
\label{eq:neumann_inverse}
\end{equation}
then converges in \(\mathcal A_\sigma\).  Uniqueness identifies it with
\(z_\vartheta\), and the representation theorem gives
\eqref{eq:operator_inverse_identity}.
\end{proof}

\begin{warning}[Exact border of the proved inverse]
\label{warn:inverse_border}
The M\"obius operator exists as a bounded operator for every \(\sigma>1\).
The argument above proves a bounded inverse only for
\(\sigma>\sigma_*\).  At \(\sigma=\sigma_*\) the Neumann norm equals
one and this argument is inconclusive.  No claim is made here for
\(1<\sigma\le\sigma_*\).
\end{warning}

\section{The first Heisenberg arithmetic law}
\label{sec:prime_pair_interference}

\begin{theorem}[Prime rays are rigid; distinct primes interfere]
\label{thm:prime_pair_interference}
Let $p,q$ be distinct primes and let $r\ge1$.  Then
\begin{align}
z_{\vartheta}(p^r)
&=1,
\label{eq:ztheta_prime_power}\\
z_{\vartheta}(pq)
&=
2\cos\!\left(\frac{\vartheta}{2}\log\frac{q}{p}\right)-1,
\label{eq:ztheta_semiprime_cos}\\
z_{\vartheta}(pq)-1
&=
-4\sin^2\!\left(\frac{\vartheta}{4}\log\frac{q}{p}\right).
\label{eq:ztheta_semiprime_defect}
\end{align}
\end{theorem}

\begin{proof}
For $p^r$, only the divisors $1$ and $p$ contribute to the recursion because
$\mu(p^j)=0$ for $j\ge2$.  Moreover,
\begin{equation}
\Om(p)\log(p^{r-1})
-
\log p\,\Om(p^{r-1})
=0,
\label{eq:prime_power_symplectic_zero}
\end{equation}
so $\chi_{\vartheta}(p,p^{r-1})=1$.  Equation
\eqref{eq:ztheta_recursion} gives
\begin{equation}
z_{\vartheta}(p^r)=z_{\vartheta}(p^{r-1}),
\label{eq:prime_power_recursion}
\end{equation}
and \eqref{eq:ztheta_prime_power} follows from $z_{\vartheta}(1)=1$.

For $pq$, the four divisors $1,p,q,pq$ give
\begin{equation}
0
=
z_{\vartheta}(pq)
-
\chi_{\vartheta}(p,q)
-
\chi_{\vartheta}(q,p)
+1.
\label{eq:pq_inverse_equation}
\end{equation}
Since
\begin{equation}
\chi_{\vartheta}(p,q)
=
\exp\!\left(\frac{i\vartheta}{2}\log\frac{q}{p}\right),
\label{eq:chi_pq}
\end{equation}
and $\chi_{\vartheta}(q,p)$ is its complex conjugate,
\eqref{eq:ztheta_semiprime_cos} follows.  The identity
\eqref{eq:ztheta_semiprime_defect} is the elementary formula
$2(\cos u-1)=-4\sin^2(u/2)$.
\end{proof}

\begin{corollary}[Loss of ordinary multiplicativity]
\label{cor:loss_of_multiplicativity}
Generically,
\begin{equation}
z_{\vartheta}(pq)
\neq
z_{\vartheta}(p)z_{\vartheta}(q)=1.
\label{eq:loss_of_multiplicativity}
\end{equation}
Thus the Heisenberg deformation preserves every one-prime ray but destroys
ordinary multiplicativity through cross-prime interference.
\end{corollary}

\begin{corollary}[Small-area expansion]
\label{cor:small_area_expansion}
For fixed distinct primes $p,q$,
\begin{equation}
z_{\vartheta}(pq)
=
1-
\frac{\vartheta^2}{4}
\left(\log\frac{q}{p}\right)^2
+O(\vartheta^4)
\qquad(\vartheta\to0).
\label{eq:small_theta_pq}
\end{equation}
There is no linear correction.
\end{corollary}

\begin{proposition}[Two-prime lattice recurrence]
\label{prop:two_prime_recurrence}
Let \(p\ne q\) be primes, set
\begin{equation}
Z_{a,b}=z_\vartheta(p^aq^b),
\qquad
\lambda=\frac{\vartheta}{2}\log\frac{q}{p}.
\label{eq:Zab_lambda}
\end{equation}
Then for \(a,b\ge1\),
\begin{equation}
Z_{a,b}
=
e^{ib\lambda}Z_{a-1,b}
+e^{-ia\lambda}Z_{a,b-1}
-e^{i(b-a)\lambda}Z_{a-1,b-1},
\label{eq:two_prime_recurrence}
\end{equation}
with boundary values
\begin{equation}
Z_{a,0}=Z_{0,b}=1.
\label{eq:two_prime_boundary}
\end{equation}
In particular,
\begin{equation}
z_\vartheta(p^2q)
=1+2\cos(2\lambda)-2\cos\lambda.
\label{eq:p2q_formula}
\end{equation}
\end{proposition}

\begin{proof}
Apply the right-inverse recursion \eqref{eq:ztheta_recursion} to
\(p^aq^b\).  Since \(\mu(d)\neq0\) only when the exponents of \(p,q\)
in \(d\) are at most one, only \(d=p,q,pq\) contribute.  Evaluation of
the three cocycle factors gives \eqref{eq:two_prime_recurrence}; the boundary
values are Theorem~\ref{thm:prime_pair_interference} on prime rays.
\end{proof}

\begin{proposition}[Evenness of the full inverse]
\label{prop:theta_evenness}
For every $n\ge1$,
\begin{equation}
z_{-\vartheta}(n)=z_{\vartheta}(n).
\label{eq:ztheta_evenness}
\end{equation}
\end{proposition}

\begin{proof}
Changing $\vartheta$ to $-\vartheta$ reverses the cocycle:
\begin{equation}
(f\star_{-\vartheta}g)(n)
=
(g\star_{\vartheta}f)(n).
\label{eq:opposite_algebra_identity}
\end{equation}
Since $z_{\vartheta}$ is a two-sided inverse of $\mu$ for
$\star_{\vartheta}$, it is a two-sided inverse of $\mu$ for the opposite
product $\star_{-\vartheta}$.  Uniqueness of the inverse gives
\eqref{eq:ztheta_evenness}.
\end{proof}

\begin{proposition}[Reality for real deformation]
\label{prop:ztheta_real}
For real \(\vartheta\) and every \(n\ge1\),
\begin{equation}
z_\vartheta(n)\in\mathbb R.
\label{eq:ztheta_real}
\end{equation}
\end{proposition}

\begin{proof}
Complex conjugation changes \(\chi_\vartheta\) to
\(\chi_{-\vartheta}\).  Since \(\mu\) is real, conjugating the inverse
equation shows that \(\overline{z_\vartheta}\) is the inverse for the
\(-\vartheta\) product.  Proposition~\ref{prop:theta_evenness} therefore
gives \(\overline{z_\vartheta}=z_{-\vartheta}=z_\vartheta\).
\end{proof}

\begin{remark}[Why this is the non-tautological step]
\label{rem:non_tautological_step}
At $\vartheta=0$ the inverse coefficients are identically $1$.  At nonzero
$\vartheta$, equation~\eqref{eq:ztheta_semiprime_defect} creates a new
arithmetic coefficient depending on the logarithmic separation of two
primes.  This dependence is forced by the symplectic area in the Weyl
product and disappears if the CCR is removed.  It is therefore the first
result in the notebook where Heisenberg noncommutativity changes the
arithmetic inverse itself rather than merely constraining an unrelated
state space.
\end{remark}

\section{Cohomological curvature of the twisted M\"obius inverse}
\label{sec:cohomological_curvature}

Set
\begin{equation}
\mathcal A(d,e)=\Omega(d)\log e-\log d\,\Omega(e),
\label{eq:arith_symplectic_area_curv}
\end{equation}
so that
\begin{equation}
\chi_\vartheta(d,e)=\exp\!\left(\frac{i\vartheta}{2}\mathcal A(d,e)\right).
\label{eq:chi_area_curv}
\end{equation}

\begin{proposition}[Coefficientwise entire dependence]
\label{prop:ztheta_entire_coefficientwise}
For each fixed \(n\ge1\), the map \(\vartheta\mapsto z_\vartheta(n)\) is
entire. It is even, and it is real-valued on the real axis.
\end{proposition}

\begin{proof}
Formula~\eqref{eq:ordered_factorization_formula} is a finite sum for fixed
\(n\), each term being an exponential \(e^{ic\vartheta}\). Evenness and
reality follow from Propositions~\ref{prop:theta_evenness} and
\ref{prop:ztheta_real}.
\end{proof}

\begin{proposition}[Second derivative of the inverse]
\label{prop:second_derivative_inverse}
Define
\begin{equation}
g(k)=\frac14\sum_{d\mid k}\mu(d)\mathcal A(d,k/d)^2.
\label{eq:gk_definition}
\end{equation}
Then
\begin{equation}
\boxed{z_0''(n)=\sum_{k\mid n}g(k).}
\label{eq:zsecond_divisor_sum}
\end{equation}
\end{proposition}

\begin{proof}
Differentiate the inverse equation
\begin{equation}
0=\sum_{d\mid n}\mu(d)z_\vartheta(n/d)
\exp\!\left(\frac{i\vartheta}{2}\mathcal A(d,n/d)\right)
\label{eq:inverse_equation_theta}
\end{equation}
twice at \(\vartheta=0\). Since \(z_0\equiv1\) and \(z'_0\equiv0\),
\begin{equation}
(\mu*z_0'')(n)=\frac14\sum_{d\mid n}\mu(d)\mathcal A(d,n/d)^2=g(n).
\label{eq:mu_convolution_zsecond}
\end{equation}
Convolving with \(\mathbf1=\mu^{-1}\) gives
\eqref{eq:zsecond_divisor_sum}.
\end{proof}

\begin{lemma}[Support collapse]
\label{lem:gk_support_collapse}
Let \(k=\prod_{j=1}^{r}p_j^{a_j}\) with distinct primes \(p_j\). Then
\begin{equation}
g(k)=0\qquad\text{for }r=0,\ r=1,\text{ or }r\ge3.
\label{eq:gk_support_zero}
\end{equation}
For two distinct primes,
\begin{equation}
\boxed{g(p^\alpha q^\beta)
=-\frac{\alpha\beta}{2}\log^2\!\frac qp.}
\label{eq:g_two_prime_exact}
\end{equation}
\end{lemma}

\begin{proof}
A contributing divisor is square-free and therefore corresponds to a subset
\(S\subseteq\{1,\ldots,r\}\). Put
\begin{equation}
R=\Omega(k)=\sum_j a_j,\qquad L=\log k,
\qquad X_S=\sum_{j\in S}\log p_j.
\label{eq:R_L_XS}
\end{equation}
Then
\begin{equation}
\mathcal A(d,k/d)=|S|L-RX_S.
\label{eq:A_subset_linear}
\end{equation}
Its square has Boolean degree at most two, so its full alternating difference
vanishes when \(r\ge3\). For \(r=1\) the area is zero. For \(r=2\), direct
evaluation of the four subsets gives
\begin{equation}
\sum_{d\mid p^\alpha q^\beta}\mu(d)
\mathcal A(d,p^\alpha q^\beta/d)^2
=-2\alpha\beta\log^2(q/p),
\label{eq:four_term_two_prime_sum}
\end{equation}
which yields \eqref{eq:g_two_prime_exact}.
\end{proof}

\begin{theorem}[Exact cohomological curvature]
\label{thm:exact_curvature_formula}
If
\begin{equation}
n=\prod_{j=1}^{r}p_j^{a_j},\qquad T(a)=\frac{a(a+1)}2,
\label{eq:n_factorization_triangular}
\end{equation}
then
\begin{equation}
\boxed{
z_0''(n)=-\frac12\sum_{1\le i<j\le r}
T(a_i)T(a_j)\log^2\!\frac{p_j}{p_i}.}
\label{eq:zsecond_general_formula}
\end{equation}
Thus \(-z_0''(n)\ge0\), with equality exactly for \(n=1\) or a prime
power.
\end{theorem}

\begin{proof}
Only divisors supported on exactly two primes survive by
Lemma~\ref{lem:gk_support_collapse}. Summing
\eqref{eq:g_two_prime_exact} over
\(1\le\alpha\le a_i\), \(1\le\beta\le a_j\) gives the factor
\(T(a_i)T(a_j)\).
\end{proof}

\begin{corollary}[Square-free prime-log variance]
\label{cor:squarefree_prime_variance}
If \(n=p_1\cdots p_r\) is square-free, then
\begin{align}
-z_0''(n)
&=\frac12\sum_{i<j}(\log p_j-\log p_i)^2
\label{eq:squarefree_pair_variance}\\
&=\frac12\left[r\sum_{p\mid n}(\log p)^2-(\log n)^2\right]
\label{eq:squarefree_variance_identity}\\
&=\frac{r^2}{2}\operatorname{Var}_{p\mid n}(\log p).
\label{eq:squarefree_variance_stat}
\end{align}
\end{corollary}

\begin{corollary}[Infinitesimal cohomological size]
\label{cor:cohomological_singular_value}
For square-free \(n=p_1\cdots p_r\), let
\begin{equation}
B_{ij}=\log p_j-\log p_i.
\label{eq:B_matrix_logs}
\end{equation}
Then \(B\) has rank at most two. If \(r\ge2\), its two nonzero singular
values are both
\begin{equation}
s(n)=\sqrt{r\sum_{p\mid n}(\log p)^2-(\log n)^2},
\label{eq:singular_value_s}
\end{equation}
and
\begin{equation}
\boxed{-z_0''(n)=\frac12s(n)^2.}
\label{eq:curvature_singular_value}
\end{equation}
\end{corollary}

\begin{proof}
Write \(B=\mathbf1a^T-a\mathbf1^T\), \(a_i=\log p_i\). The two nonzero
singular values of \(uv^T-vu^T\) equal
\begin{equation}
\sqrt{\|u\|^2\|v\|^2-(u\cdot v)^2}.
\label{eq:wedge_singular_value}
\end{equation}
Apply this with \(u=\mathbf1\), \(v=a\), and then use
Corollary~\ref{cor:squarefree_prime_variance}.
\end{proof}

\begin{corollary}[Exact Taylor law for each integer]
\label{cor:coefficient_taylor_exact}
For every fixed \(n\),
\begin{equation}
z_\vartheta(n)=1+\frac{\vartheta^2}{2}z_0''(n)+O_n(\vartheta^4),
\label{eq:ztheta_general_taylor}
\end{equation}
and the Taylor series has infinite radius of convergence. For square-free
\(n\),
\begin{equation}
\boxed{
z_\vartheta(n)=1-\frac{\vartheta^2}{4}
\left[\omega(n)\sum_{p\mid n}(\log p)^2-(\log n)^2\right]
+O_n(\vartheta^4).}
\label{eq:ztheta_squarefree_taylor}
\end{equation}
\end{corollary}

\section{Dirichlet summation of the curvature}
\label{sec:curvature_dirichlet_series}

Define
\begin{equation}
c(n)=-z_0''(n)\ge0.
\label{eq:cn_curvature}
\end{equation}

\begin{theorem}[Exact curvature Dirichlet series]
\label{thm:curvature_dirichlet_series}
For \(\Re s>1\),
\begin{equation}
\mathcal K(s)=\sum_{n\ge1}\frac{c(n)}{n^s}
\label{eq:K_dirichlet_definition}
\end{equation}
converges absolutely and
\begin{equation}
\boxed{
\mathcal K(s)=\frac{\zeta(s)}2\left[S_0(s)S_2(s)-S_1(s)^2\right],}
\label{eq:K_variance_determinant}
\end{equation}
where
\begin{equation}
S_j(s)=\sum_p\frac{p^{-s}}{(1-p^{-s})^2}(\log p)^j,
\qquad j=0,1,2.
\label{eq:Sj_prime_sum}
\end{equation}
Equivalently,
\begin{equation}
S_0S_2-S_1^2
=\sum_{p<q}
\frac{p^{-s}q^{-s}}{(1-p^{-s})^2(1-q^{-s})^2}
\log^2\!\frac qp.
\label{eq:weighted_prime_variance}
\end{equation}
\end{theorem}

\begin{proof}
Use Theorem~\ref{thm:exact_curvature_formula} and
\begin{equation}
\sum_{a\ge1}T(a)x^a=\frac{x}{(1-x)^3}.
\label{eq:triangular_generating_function}
\end{equation}
Factoring the unrestricted Euler factors gives \(\zeta(s)\), while a selected
prime contributes
\begin{equation}
\frac{p^{-s}}{(1-p^{-s})^2}.
\label{eq:selected_prime_weight}
\end{equation}
The pair sum is exactly the weighted variance determinant
\(S_0S_2-S_1^2\). Absolute convergence follows for \(\Re s>1\).
\end{proof}

\begin{corollary}[Prime-zeta reduction]
\label{cor:prime_zeta_representation}
Let \(P(s)=\sum_p p^{-s}\) for \(\Re s>1\). Then
\begin{align}
S_0(s)&=\sum_{m\ge1}mP(ms),
\label{eq:S0_prime_zeta}\\
S_1(s)&=-\sum_{m\ge1}mP'(ms),
\label{eq:S1_prime_zeta}\\
S_2(s)&=\sum_{m\ge1}mP''(ms).
\label{eq:S2_prime_zeta}
\end{align}
Using
\begin{equation}
P(s)=\sum_{k\ge1}\frac{\mu(k)}{k}\log\zeta(ks),
\label{eq:prime_zeta_mobius_inversion}
\end{equation}
gives
\begin{align}
S_0(s)&=\sum_{r\ge1}\frac{J_2(r)}{r}\log\zeta(rs),
\label{eq:S0_logzeta}\\
S_1(s)&=-\sum_{r\ge1}\varphi(r)\frac{\zeta'(rs)}{\zeta(rs)},
\label{eq:S1_logderivative}\\
S_2(s)&=\left(\frac{\zeta'}{\zeta}\right)'(s),
\label{eq:S2_logderivative_exact}
\end{align}
where \(J_2(r)=r^2\prod_{p\mid r}(1-p^{-2})\).
\end{corollary}

\begin{proof}
Expand \(x/(1-x)^2=\sum_{m\ge1}mx^m\), substitute
\eqref{eq:prime_zeta_mobius_inversion}, and group by \(r=mk\). The needed
divisor identities are
\begin{align}
r\sum_{k\mid r}\frac{\mu(k)}{k^2}&=\frac{J_2(r)}r,
\label{eq:J2_divisor_identity}\\
r\sum_{k\mid r}\frac{\mu(k)}k&=\varphi(r),
\label{eq:phi_divisor_identity}\\
r\sum_{k\mid r}\mu(k)&=\begin{cases}1,&r=1,\\0,&r>1.\end{cases}
\label{eq:mobius_divisor_delta}
\end{align}
The last identity collapses \(S_2\) to
\eqref{eq:S2_logderivative_exact}.
\end{proof}

\begin{proposition}[Singularity at \(s=1\)]
\label{prop:K_near_one}
As \(s\to1^+\),
\begin{equation}
\boxed{
\mathcal K(s)\sim\frac1{2(s-1)^3}\log\frac1{s-1}.}
\label{eq:K_near_one_asymptotic}
\end{equation}
\end{proposition}

\begin{proof}
From \eqref{eq:prime_zeta_mobius_inversion},
\begin{equation}
P(s)=\log\frac1{s-1}+O(1).
\label{eq:P_near_one}
\end{equation}
Hence
\begin{equation}
S_0(s)=\log\frac1{s-1}+O(1),\quad
S_1(s)=\frac1{s-1}+O(1),\quad
S_2(s)=\frac1{(s-1)^2}+O(1).
\label{eq:S012_near_one}
\end{equation}
Combine with \(\zeta(s)\sim(s-1)^{-1}\).
\end{proof}

\section{Continuation of the curvature and detection of zeta zeros}
\label{sec:curvature_zero_detection}

\begin{theorem}[Local continuation to \(\Re s>1/2\)]
\label{thm:K_local_continuation}
Let \(U\subset\{\Re s>1/2\}\) be simply connected and contain no zero or pole
of \(\zeta(s)\). Choose a branch of \(\log\zeta(s)\) on \(U\). Then
\eqref{eq:S0_logzeta}--\eqref{eq:S2_logderivative_exact}, substituted into
\eqref{eq:K_variance_determinant}, continue \(\mathcal K\) holomorphically to
\(U\).
\end{theorem}

\begin{proof}
For \(r\ge2\), \(\Re(rs)>1\), so the corresponding terms are holomorphic and
converge normally on compact subsets. The \(r=1\) terms are holomorphic by the
hypothesis on \(U\).
\end{proof}

\begin{theorem}[Zeta zeros are logarithmic curvature singularities]
\label{thm:zero_curvature_singularity}
Let \(\rho\) be a zero of multiplicity \(m\ge1\) with
\(1/2<\Re\rho<1\), and write
\begin{equation}
\zeta(s)=c_\rho(s-\rho)^m(1+O(s-\rho)),\qquad c_\rho\ne0.
\label{eq:zeta_local_zero_m}
\end{equation}
Then
\begin{equation}
\boxed{
\mathcal K(s)=-\frac{c_\rho m^2}{2}
(s-\rho)^{m-2}\log(s-\rho)
+O\!\left((s-\rho)^{m-2}\right).}
\label{eq:K_zero_local_form}
\end{equation}
For a simple zero,
\begin{equation}
\boxed{
\mathcal K(s)=-\frac{\zeta'(\rho)}2
\frac{\log(s-\rho)}{s-\rho}
+O\!\left(\frac1{s-\rho}\right).}
\label{eq:K_simple_zero_form}
\end{equation}
\end{theorem}

\begin{proof}
For \(r\ge2\), the terms in
\eqref{eq:S0_logzeta}--\eqref{eq:S1_logderivative} are holomorphic near
\(\rho\). With \(t=s-\rho\), the singular pieces are
\begin{equation}
S_0=m\log t+O(1),\qquad
S_1=-\frac mt+O(1),\qquad
S_2=-\frac m{t^2}+O(1).
\label{eq:S012_near_zero}
\end{equation}
Thus
\begin{equation}
S_0S_2-S_1^2=-\frac{m^2\log t}{t^2}+O(t^{-2}).
\label{eq:det_near_zero}
\end{equation}
Multiply by \(\zeta(s)/2\).
\end{proof}

\begin{corollary}[Curvature criterion equivalent to RH]
\label{cor:curvature_RH_criterion}
RH is equivalent to the absence of logarithmic branch singularities of
\(\mathcal K\) in
\begin{equation}
\boxed{\frac12<\Re s<1.}
\label{eq:RH_curvature_halfstrip}
\end{equation}
\end{corollary}

\begin{proof}
Theorem~\ref{thm:zero_curvature_singularity} detects every zero in the right
half of the critical strip, while Theorem~\ref{thm:K_local_continuation}
gives holomorphy on zero-free neighbourhoods. The functional equation makes
the absence of zeros in this half-strip equivalent to RH.
\end{proof}

\section{The Heisenberg-deformed zeta series}
\label{sec:deformed_zeta_series}

\begin{definition}[Heisenberg-deformed zeta series]
\label{def:deformed_zeta_series}
For $\Re s>\sigma_*$, define
\begin{equation}
\zeta_{\mathrm H}(s;\vartheta)
=
\sum_{n\ge1}\frac{z_{\vartheta}(n)}{n^s}.
\label{eq:deformed_zeta_series}
\end{equation}
Theorem~\ref{thm:neumann_region} guarantees absolute convergence for every $\Re s>\sigma_*$.
\end{definition}

\begin{proposition}[Classical boundary]
\label{prop:zeta_classical_boundary}
For $\Re s>1$,
\begin{equation}
\zeta_{\mathrm H}(s;0)=\zeta(s).
\label{eq:zeta_classical_boundary}
\end{equation}
\end{proposition}

\begin{proof}
By Proposition~\ref{prop:commuting_boundary}, $z_0(n)=1$ for all $n$.
Substitution into \eqref{eq:deformed_zeta_series} gives the ordinary
Dirichlet series for $\zeta(s)$.
\end{proof}

\begin{proposition}[Even Heisenberg deformation]
\label{prop:zeta_theta_even}
For $\Re s>\sigma_*$,
\begin{equation}
\zeta_{\mathrm H}(s;-\vartheta)
=
\zeta_{\mathrm H}(s;\vartheta).
\label{eq:zeta_theta_even}
\end{equation}
\end{proposition}

\begin{proof}
This follows termwise from Proposition~\ref{prop:theta_evenness} and absolute
convergence.
\end{proof}

\begin{warning}[No analytic continuation is asserted]
\label{warn:no_continuation}
The series \eqref{eq:deformed_zeta_series} is absolutely convergent and holomorphic on the open half-plane $\Re s>\sigma_*$.  No meromorphic continuation, functional equation, scalar Euler product, or critical-strip theory
for $\zeta_{\mathrm H}(s;\vartheta)$ is asserted here.
\end{warning}

\begin{proposition}[Conditional first-order rigidity of simple zeros]
\label{prop:conditional_zero_rigidity}
Assume that, for some simple zero $\rho$ of the ordinary zeta function,
there exists a function $F(s,\vartheta)$ analytic in a neighbourhood of
$(\rho,0)$ such that
\begin{align}
F(s,0)
&=\zeta(s),
\label{eq:F_classical_boundary}\\
F(s,-\vartheta)
&=F(s,\vartheta),
\label{eq:F_evenness}
\end{align}
and that $F$ agrees with $\zeta_{\mathrm H}$ wherever both are defined.
Then the local zero branch $\rho(\vartheta)$ through $\rho$ satisfies
\begin{equation}
\rho'(0)=0.
\label{eq:zero_first_derivative_zero}
\end{equation}
\end{proposition}

\begin{proof}
The implicit-function theorem applies because $\zeta'(\rho)\neq0$.
Evenness \eqref{eq:F_evenness} makes the unique local zero branch even, so
its first derivative at the origin vanishes.
\end{proof}

\begin{warning}
\label{warn:conditional_zero_result}
Proposition~\ref{prop:conditional_zero_rigidity} is conditional on an analytic
continuation in both variables and does not supply such a continuation.  It
is included because it identifies a precise way in which the exact
$\vartheta\mapsto-\vartheta$ symmetry would constrain the motion of classical
Riemann zeros under the Heisenberg deformation if that continuation exists.
\end{warning}

\section{A genuine M\"obius Hamiltonian}
\label{sec:dynamics}

\subsection{Definition and domain}
\label{subsec:hamiltonian_definition}

\begin{definition}[Cosine and sine Weyl quadratures]
\label{def:Weyl_quadratures}
For $n\ge1$, set
\begin{align}
C_n
&=\frac{W_n+W_n^*}{2},
\label{eq:Cn_definition}\\
J_n
&=\frac{W_n-W_n^*}{2i}.
\label{eq:Jn_definition}
\end{align}
Both are bounded self-adjoint operators with norm at most $1$.
\end{definition}

\begin{definition}[M\"obius Hamiltonian]
\label{def:mobius_hamiltonian}
Fix an energy scale $E_0>0$ and $\sigma>1$.  Define
\begin{equation}
H_{\mu,\sigma,\vartheta}
=
E_0\left(\Reop\mathcal M_{\sigma,\vartheta}-I\right)
=
E_0\sum_{n\ge2}
\frac{\mu(n)}{n^{\sigma}}C_n.
\label{eq:mobius_hamiltonian}
\end{equation}
The series converges absolutely in operator norm.
\end{definition}

\begin{proposition}[Self-adjoint bounded dynamics]
\label{prop:hamiltonian_selfadjoint}
For $\sigma>1$, $H_{\mu,\sigma,\vartheta}$ is bounded and self-adjoint on
all of $\Hil$:
\begin{equation}
\Dom(H_{\mu,\sigma,\vartheta})=\Hil.
\label{eq:hamiltonian_domain}
\end{equation}
It therefore generates the unitary group
\begin{equation}
U(t)=\exp\!\left(-\frac{i}{\hbar}tH_{\mu,\sigma,\vartheta}\right).
\label{eq:time_evolution_unitary}
\end{equation}
\end{proposition}

\begin{proof}
The norm series in \eqref{eq:mobius_hamiltonian} is dominated by
$E_0\sum_{n\ge2}|\mu(n)|n^{-\sigma}$, which converges for $\sigma>1$.
Every $C_n$ is self-adjoint, hence so is the norm limit.  The spectral theorem
then defines \eqref{eq:time_evolution_unitary}.
\end{proof}

\begin{proposition}[Explicit Schr\"odinger equation]
\label{prop:schrodinger_equation}
For $\psi_0\in\Hil$, the evolution $\psi(t)=U(t)\psi_0$ satisfies
\begin{equation}
i\hbar\,\partial_t\psi(t,q)
=
(H_{\mu,\sigma,\vartheta}\psi(t))(q)
\label{eq:schrodinger_equation}
\end{equation}
in the strong $L^2$ sense, where
\begin{align}
(H_{\mu,\sigma,\vartheta}\psi)(q)
=
\frac{E_0}{2}
\sum_{n\ge2}
\frac{\mu(n)}{n^{\sigma}}
\Bigg[
& e^{\frac{i}{\hbar}y_n(q-x_n/2)}\psi(q-x_n)
\notag\\
+{}&e^{-\frac{i}{\hbar}y_n(q+x_n/2)}\psi(q+x_n)
\Bigg].
\label{eq:hamiltonian_coordinate_action}
\end{align}
\end{proposition}

\begin{proof}
Equation~\eqref{eq:hamiltonian_coordinate_action} follows from
Proposition~\ref{prop:weyl_coordinate_action} applied to $W_n$ and $W_n^*$.
The series converges in operator norm, so the bounded Hamiltonian gives a
strong Schr\"odinger evolution for every initial vector.
\end{proof}

\subsection{Exact equations of motion for position and momentum}
\label{subsec:equations_of_motion}

\begin{lemma}[Bounded first commutators]
\label{lem:bounded_first_commutators}
For $\sigma>1$, the series
\begin{align}
\sum_{n\ge2}
\frac{|\mu(n)|x_n}{n^{\sigma}},
\label{eq:x_commutator_sum}\\
\sum_{n\ge2}
\frac{|\mu(n)|y_n}{n^{\sigma}}
\label{eq:y_commutator_sum}
\end{align}
converge.  Hence $[H_{\mu,\sigma,\vartheta},Q]$ and
$[H_{\mu,\sigma,\vartheta},P]$, initially defined on $\Sch(\R)$, extend to
bounded operators.
\end{lemma}

\begin{proof}
The first series is bounded by a constant multiple of
$\sum n^{-\sigma}\log n$.  Since
\begin{equation}
\Om(n)\le\frac{\log n}{\log2},
\label{eq:Omega_log_bound}
\end{equation}
the second is bounded by another constant multiple of the same convergent
series.
\end{proof}

\begin{theorem}[Exact Heisenberg equations of motion]
\label{thm:equations_of_motion}
Let
\begin{align}
Q(t)
&=U(t)^*QU(t),
\label{eq:Qt_definition}\\
P(t)
&=U(t)^*PU(t).
\label{eq:Pt_definition}
\end{align}
The bounded commutators in Lemma~\ref{lem:bounded_first_commutators} imply that $U(t)$ preserves $\Dom(Q)$ and $\Dom(P)$.  On these domains, and equivalently as quadratic-form identities, the Heisenberg equations are
\begin{align}
\dot Q(t)
&=
\frac{E_0}{\hbar}
U(t)^*
\left[
\sum_{n\ge2}
\frac{\mu(n)x_n}{n^{\sigma}}J_n
\right]
U(t),
\label{eq:Q_equation_motion}\\
\dot P(t)
&=
\frac{E_0}{\hbar}
U(t)^*
\left[
\sum_{n\ge2}
\frac{\mu(n)y_n}{n^{\sigma}}J_n
\right]
U(t).
\label{eq:P_equation_motion}
\end{align}
At $t=0$ this becomes
\begin{align}
\dot Q(0)
&=
\frac{E_0\ell}{\hbar}
\sum_{n\ge2}
\frac{\mu(n)\log n}{n^{\sigma}}J_n,
\label{eq:Qdot_zero}\\
\dot P(0)
&=
\frac{E_0\kappa}{\hbar}
\sum_{n\ge2}
\frac{\mu(n)\Om(n)}{n^{\sigma}}J_n.
\label{eq:Pdot_zero}
\end{align}
\end{theorem}

\begin{proof}
For $A=Q$ or $P$, the bounded-commutator theorem for unitary dynamics gives
\begin{equation}
\dot A(t)
=
\frac{i}{\hbar}U(t)^*[H_{\mu,\sigma,\vartheta},A]U(t).
\label{eq:general_heisenberg_equation}
\end{equation}
From Proposition~\ref{prop:weyl_QP_commutators},
\begin{align}
[Q,C_n]
&=ix_nJ_n,
\label{eq:Q_Cn_commutator}\\
[P,C_n]
&=iy_nJ_n.
\label{eq:P_Cn_commutator}
\end{align}
Therefore
\begin{align}
[H_{\mu,\sigma,\vartheta},Q]
&=
-iE_0\sum_{n\ge2}
\frac{\mu(n)x_n}{n^{\sigma}}J_n,
\label{eq:H_Q_commutator}\\
[H_{\mu,\sigma,\vartheta},P]
&=
-iE_0\sum_{n\ge2}
\frac{\mu(n)y_n}{n^{\sigma}}J_n.
\label{eq:H_P_commutator}
\end{align}
Lemma~\ref{lem:bounded_first_commutators} justifies the sums as bounded
operators.  Substitution into \eqref{eq:general_heisenberg_equation} gives
\eqref{eq:Q_equation_motion} and \eqref{eq:P_equation_motion}.
\end{proof}

\begin{remark}[Physical and arithmetic meaning of the two equations]
\label{rem:meaning_equations_motion}
The position variable is $Q$, the logarithmic-scale coordinate.  The
momentum variable is $P=-i\hbar\partial_q$, the generator conjugate to that
coordinate.  Multiplication by an integer $n$ is represented by a phase-space
displacement of size $\ell\log n$ in $Q$ and $\kappa\Om(n)$ in $P$.
Equations~\eqref{eq:Qdot_zero} and \eqref{eq:Pdot_zero} therefore say that
the M\"obius Hamiltonian drives position by logarithmic size and momentum by
prime-factor count, with the same quantum current operators $J_n$.
\end{remark}

\subsection{The CCR is visible in the dynamics}
\label{subsec:CCR_in_dynamics}

\begin{proposition}[Pairwise dynamical noncommutativity]
\label{prop:CJ_commutator}
Extend $\log$ and $\Om$ additively to positive rationals and define $W_r$ by
\eqref{eq:Wn_definition} for $r\in\Qp$.  For $m,n\in\N$, put
\begin{equation}
\Phi_{m,n}
=
\vartheta\bigl(\Om(m)\log n-\log m\,\Om(n)\bigr).
\label{eq:Phi_mn}
\end{equation}
Then
\begin{align}
[C_m,J_n]
=
\frac12\sin\!\left(\frac{\Phi_{m,n}}{2}\right)
\Bigl(
W_{mn}+W_{m/n}-W_{n/m}-W_{1/(mn)}
\Bigr).
\label{eq:Cm_Jn_commutator}
\end{align}
\end{proposition}

\begin{proof}
Expand $C_m$ and $J_n$ using
\eqref{eq:Cn_definition}--\eqref{eq:Jn_definition}.  Apply
\eqref{eq:weyl_commutator} to the four pairs
$(z_m,z_n)$, $(z_m,-z_n)$, $(-z_m,z_n)$, and $(-z_m,-z_n)$.
\end{proof}

\begin{corollary}[Second-order arithmetic commutator]
\label{cor:second_order_force}
At $t=0$,
\begin{align}
\ddot Q(0)
=
\frac{iE_0^2}{2\hbar^2}
\sum_{m,n\ge2}
&\frac{\mu(m)\mu(n)x_n}{(mn)^{\sigma}}
\sin\!\left(\frac{\Phi_{m,n}}{2}\right)
\notag\\
&\times
\Bigl(
W_{mn}+W_{m/n}-W_{n/m}-W_{1/(mn)}
\Bigr),
\label{eq:Q_second_derivative}
\end{align}
and the analogous momentum equation is
\begin{align}
\ddot P(0)
=
\frac{iE_0^2}{2\hbar^2}
\sum_{m,n\ge2}
&\frac{\mu(m)\mu(n)y_n}{(mn)^{\sigma}}
\sin\!\left(\frac{\Phi_{m,n}}{2}\right)
\notag\\
&\times
\Bigl(
W_{mn}+W_{m/n}-W_{n/m}-W_{1/(mn)}
\Bigr).
\label{eq:P_second_derivative}
\end{align}
The double sums converge in operator norm for $\sigma>1$.
\end{corollary}

\begin{proof}
Differentiate \eqref{eq:Q_equation_motion} and
\eqref{eq:P_equation_motion} once more and use
Proposition~\ref{prop:CJ_commutator}.  Polynomial logarithmic weights are
summable against $(mn)^{-\sigma}$ for $\sigma>1$.
\end{proof}

\begin{corollary}[The interaction disappears with the Weyl cocycle]
\label{cor:interaction_disappears}
If $\vartheta=0$, then
\begin{equation}
\Phi_{m,n}=0
\label{eq:Phi_zero}
\end{equation}
for all $m,n$, and all pairwise terms in
\eqref{eq:Q_second_derivative}--\eqref{eq:P_second_derivative} vanish.
Thus the second-order cross-arithmetic second-order term displayed above is generated
entirely by the exact Heisenberg cocycle.
\end{corollary}

\section{Exact non-localization inside the arithmetic dynamics}
\label{sec:nonlocalization}

\begin{proposition}[Schwartz action and distributional transpose]
\label{prop:schwartz_distribution_extension}
For \(\sigma>1\), the Hamiltonian
\(H_{\mu,\sigma,\vartheta}\) maps \(\Sch(\R)\) continuously into
itself.  Its continuous transpose
\(H^{\times}:\Sch'(\R)\to\Sch'(\R)\) is defined by
\begin{equation}
\langle H^{\times}T,\phi\rangle
=
\langle T,H_{\mu,\sigma,\vartheta}\phi\rangle,
\qquad
T\in\Sch'(\R),\quad\phi\in\Sch(\R).
\label{eq:distributional_transpose_definition}
\end{equation}
Then \(H^{\times}\) is continuous on \(\Sch'(\R)\).
\end{proposition}

\begin{proof}
For every Schwartz seminorm
\(p_{a,b}(\phi)=\sup_q|q^a\phi^{(b)}(q)|\), the coordinate action of a
Weyl operator gives
\begin{equation}
p_{a,b}(W_n\phi)
\le
C_{a,b}(1+|x_n|+|y_n|)^{N_{a,b}}
\sum_{j\le b}p_{a+j,b-j}(\phi)
\label{eq:weyl_schwartz_bound}
\end{equation}
for suitable constants.  Since \(x_n,y_n=O(\log n)\) and
\begin{equation}
\sum_{n\ge2}n^{-\sigma}(1+\log n)^N<\infty
\qquad(\sigma>1)
\label{eq:log_weight_summability}
\end{equation}
for every fixed \(N\), the Hamiltonian series converges in every Schwartz
seminorm and defines a continuous map \(\Sch\to\Sch\).  Therefore
\eqref{eq:distributional_transpose_definition} defines the standard
continuous dual action on \(\Sch'\).  No termwise extension to all tempered
distributions is assumed without this transpose construction.
\end{proof}

\begin{proposition}[No exact position point is a stationary generalized mode]
\label{prop:no_delta_eigenstate}
Let \(\sigma>1\), \(\ell>0\), and \(E_0>0\).  For no
\(q_0\in\R\) and no scalar \(E\in\C\) does
\(\delta_{q_0}\in\Sch'(\R)\) satisfy
\begin{equation}
H^{\times}_{\mu,\sigma,\vartheta}\delta_{q_0}
=E\delta_{q_0}.
\label{eq:delta_eigenstate_impossible}
\end{equation}
\end{proposition}

\begin{proof}
The transpose action of a Weyl displacement on a point mass is
\begin{align}
W(x,y)^{\times}\delta_{q_0}
&=
\exp\left(\frac{i}{\hbar}y(q_0+x/2)\right)
\delta_{q_0+x},
\label{eq:W_delta_action}\\
(W(x,y)^*)^{\times}\delta_{q_0}
&=
\exp\left(-\frac{i}{\hbar}y(q_0-x/2)\right)
\delta_{q_0-x}.
\label{eq:Wstar_delta_action}
\end{align}
The \(n=2\) term therefore contributes nonzero masses at
\(q_0\pm\ell\log2\).  No other positive integer gives the positive
shift \(\ell\log2\), and no adjoint term gives that same positive shift.
Since \(\mu(2)=-1\), the corresponding coefficient cannot cancel.  Hence
the support of the left side of \eqref{eq:delta_eigenstate_impossible}
contains a point distinct from \(q_0\), while the right side is supported
only at \(q_0\).
\end{proof}

\begin{remark}[Relation to exact Heisenberg impossibility]
\label{rem:nonlocalization_relation}
Theorem~\ref{thm:no_joint_eigenstate} is a universal consequence of the CCR.
Proposition~\ref{prop:no_delta_eigenstate} is more specific: the M\"obius
Hamiltonian itself does not admit an exactly position-localized stationary
mode.  The reason is not a measurement uncertainty but the exact translation
content of the arithmetic Weyl operators.
\end{remark}

\section{Status classification in Version 4}
\label{sec:status_classification}

\subsection{Standard ingredients}
\label{subsec:standard_ingredients}

The following ingredients are standard and are not presented as new:

\begin{enumerate}[leftmargin=*]
\item the Schr\"odinger representation of $Q$ and $P$ and the CCR
\eqref{eq:CCR};
\item Weyl operators and the Weyl product
\eqref{eq:weyl_product};
\item the arithmetic functions $\mu$, $\Om$, and classical M\"obius
inversion;
\item weighted $\ell^1$ Banach-algebra estimates and Neumann inversion.
\end{enumerate}

\subsection{Original notebook constructions}

Version 4 adds:
\begin{enumerate}[label=(V4.\arabic*)]
\item the quadratic-obstruction theorem
Theorem~\ref{thm:linear_state_quadratic_obstruction};
\item coefficientwise entire dependence of \(z_\vartheta(n)\);
\item the exact curvature law \eqref{eq:zsecond_general_formula};
\item the square-free prime-log variance identity
\eqref{eq:squarefree_variance_stat};
\item the curvature Dirichlet determinant
\eqref{eq:K_variance_determinant};
\item the identity \(S_2=(\zeta'/\zeta)'\);
\item the logarithmic zero singularity \eqref{eq:K_zero_local_form};
\item the curvature criterion Corollary~\ref{cor:curvature_RH_criterion}.
\end{enumerate}

\label{subsec:original_notebook_constructions}

The following are the constructions/results introduced in this notebook and
are the candidates for substantive originality:

\begin{enumerate}[leftmargin=*]
\item the arithmetic phase-space embedding
$n\mapsto(\ell\log n,\kappa\Om(n))$ coupled to the Weyl representation;
\item the cocycle \eqref{eq:chi_definition} as a deformation of Dirichlet
convolution;
\item the arithmetic uncertainty law \eqref{eq:arithmetic_uncertainty}, which specializes the Heisenberg--Robertson inequality to the integer generators \(G_n\);
\item the two-sided Heisenberg-twisted M\"obius inverse
$z_{\vartheta}$;
\item the exact prime-ray/prime-pair law
\eqref{eq:ztheta_prime_power}--\eqref{eq:ztheta_semiprime_defect};
\item the evenness theorem \eqref{eq:ztheta_evenness};
\item the M\"obius Hamiltonian \eqref{eq:mobius_hamiltonian} and its exact
position--momentum equations
\eqref{eq:Q_equation_motion}--\eqref{eq:P_equation_motion};
\item the pairwise second-order commutator controlled by
$\sin(\Phi_{m,n}/2)$.
\end{enumerate}


\section{Corrections and extensions carried into Version 4}
\label{sec:v2_audit_corrections}

\begin{remark}
\label{rem:v2_audit_corrections}
Version 4 incorporates the following mathematical corrections and
clarifications: the cocycle is proved nontrivial rather than merely described
as noncommutative; no faithfulness of \(\Pi_{\sigma,\vartheta}\) is
claimed; left and right inverses are constructed separately before being
identified; the uniform Neumann half-plane is sharpened to
\(\sigma>\sigma_*\); the ordered-factorization formula and a two-prime
recurrence are included; real-valuedness of \(z_\vartheta(n)\) for real
\(\vartheta\) is proved; the deformed Dirichlet series is asserted
holomorphic only on its proved open half-plane; the unitary evolution is not
claimed to preserve Schwartz space; and the action on tempered distributions
is defined by continuous transposition.
\end{remark}

\section{Open problems}
\label{sec:open_problems}

\begin{question}[Analytic continuation]
\label{q:analytic_continuation}
Does $\zeta_{\mathrm H}(s;\vartheta)$ admit a natural meromorphic continuation
from $\Re s\ge2$ to a larger domain for fixed nonzero $\vartheta$?
\end{question}

\begin{question}[Functional symmetry]
\label{q:functional_symmetry}
Is there a completion of $\zeta_{\mathrm H}(s;\vartheta)$ with a symmetry
that reduces at $\vartheta=0$ to the classical reflection $s\mapsto1-s$?
A nontrivial answer would be required before using the deformation to say
anything substantive about the critical line.
\end{question}

\begin{question}[Spectral zeros]
\label{q:spectral_zeros}
Can the zero set of a continued $\zeta_{\mathrm H}$ be realized as a spectral
or resonance set of an operator built intrinsically from
$H_{\mu,\sigma,\vartheta}$, rather than as zeros inserted into a multiplier?
\end{question}

\begin{question}[Classical-zero rigidity]
\label{q:classical_zero_rigidity}
If an analytic two-variable continuation through $\vartheta=0$ exists, does
the exact evenness in $\vartheta$ impose any second-order constraint on the
motion of the classical Riemann zeros beyond
Proposition~\ref{prop:conditional_zero_rigidity}?
\end{question}

\begin{question}[Choice of the second additive coordinate]
\label{q:additive_coordinate}
Which completely additive arithmetic functions $A(n)$, replacing $\Om(n)$
in $y_n=\kappa A(n)$, produce inequivalent deformations?  The choice
$A=\Om$ is distinguished here by its direct relation to the sign of $\mu$ on
square-free integers, but a classification is open.
\end{question}


\begin{question}[Higher cohomological curvatures]
\label{q:higher_curvatures}
Determine closed formulas for
\begin{equation}
\frac{d^{2k}}{d\vartheta^{2k}}z_\vartheta(n)\Big|_{\vartheta=0},
\qquad k\ge2,
\label{eq:higher_curvature_derivatives}
\end{equation}
and for their Dirichlet generating series. The fourth-order case is the first
test of whether the collapse to finitely many one-prime moments persists or
irreducible higher prime geometry appears.
\end{question}

\section{Notebook history}
\label{sec:notebook_history}

\begin{remark}
\label{rem:history_v1}
Version 1 used M\"obius-weighted logarithmic translations, Fourier
multipliers, an eta multiplier in the critical strip, generalized plane-wave
null modes, and a standard CCR non-localization statement.  Its main defect
was that the arithmetic multiplier and the Heisenberg pair were juxtaposed
rather than coupled.
\end{remark}

\begin{remark}
\label{rem:history_v2}
Version 2 removes the critical-strip multiplier as the central mechanism and
replaces ordinary Dirichlet convolution by the Weyl-cocycle product
\eqref{eq:twisted_convolution}.  Heisenberg noncommutativity now changes the
M\"obius inverse coefficients themselves and enters the exact equations of
motion of the M\"obius Hamiltonian.
\end{remark}


\begin{remark}
\label{rem:history_v4}
Version 4 closes the linear-state route structurally and moves the analysis to
the cocycle itself. Differentiating the twisted M\"obius inverse at
\(\vartheta=0\) produces the prime-factor curvature
\eqref{eq:zsecond_general_formula}. Summing it yields
\eqref{eq:K_variance_determinant}, whose local continuation detects every
zeta zero in \(\Re s>1/2\) through \eqref{eq:K_zero_local_form}.
\end{remark}

\begin{thebibliography}{99}

\bibitem{Heisenberg1927}
W.~Heisenberg,
\emph{\"Uber den anschaulichen Inhalt der quantentheoretischen Kinematik und
Mechanik},
Zeitschrift f\"ur Physik \textbf{43} (1927), 172--198.
doi:10.1007/BF01397280.

\bibitem{Folland1989}
G.~B.~Folland,
\emph{Harmonic Analysis in Phase Space},
Annals of Mathematics Studies 122,
Princeton University Press, Princeton, 1989.

\bibitem{ReedSimon1980}
M.~Reed and B.~Simon,
\emph{Methods of Modern Mathematical Physics, Vol.~I: Functional Analysis},
revised and enlarged ed., Academic Press, 1980.

\bibitem{Apostol1976}
T.~M.~Apostol,
\emph{Introduction to Analytic Number Theory},
Springer, New York, 1976.
doi:10.1007/978-1-4757-5579-4.

\end{thebibliography}

\end{document}
